
\documentclass{article}
\usepackage{colortbl}
\usepackage{makecell}
\usepackage{multirow}
\usepackage{supertabular}

\begin{document}

\newcounter{utterance}

\centering \large Interaction Transcript for game `referencegame', experiment `random\_grids', episode 7 with reflexion{-}Qwen3.5{-}4B.
\vspace{24pt}

{ \footnotesize  \setcounter{utterance}{1}
\setlength{\tabcolsep}{0pt}
\begin{supertabular}{c@{$\;$}|p{.15\linewidth}@{}p{.15\linewidth}p{.15\linewidth}p{.15\linewidth}p{.15\linewidth}p{.15\linewidth}}
    \# & $\;$A & \multicolumn{4}{c}{Game Master} & $\;\:$B\\
    \hline

    \theutterance \stepcounter{utterance}  
    & & \multicolumn{4}{p{0.6\linewidth}}{
        \cellcolor[rgb]{0.9,0.9,0.9}{
            \makecell[{{p{\linewidth}}}]{
                \texttt{\tiny{[P1$\langle$GM]}}
                \texttt{You have attempted this same task before and did not succeed. These are the plans you wrote after those attempts.} \\
\texttt{Trial \#1: New plan: My previous expression was too verbose and listed every single cell, which likely made it difficult to distinguish the target grid from the distractors based on unique patterns. I should have focused on identifying the most distinctive feature of the target grid that is absent in both distractors. For the target grid, the unique pattern is the "X" in the top row at the third position (row 1, col 3), whereas the distractors have "X"s in the first or second column of the top row. I will construct a concise expression that highlights this specific unique location and the surrounding cluster of X's in the middle rows, explicitly stating that this configuration distinguishes it from the other two grids.} \\
\\ 
\texttt{Trial \#2: New plan: I failed because my description was too generic and did not explicitly state the unique distinguishing feature of the target grid compared to the distractors. The target grid has a specific pattern in the top row (X at col 3, with Xs at col 2 and 4) that is not present in the same configuration in the distractors. I will construct a concise expression that specifically identifies the "X" at row 1, column 3 and notes that the cells at row 1, columns 2 and 4 are also "X", creating a horizontal triplet or specific cluster that distinguishes it from the distractors where the top row patterns differ.} \\
\\ 
\texttt{Trial \#3: New plan: I failed because my description was too detailed and failed to highlight the single most distinctive feature that separates the target grid from both distractors. The target grid is the only one with an "X" in the top row at the third column (row 1, col 3), whereas the distractors have "X"s at columns 1 or 2 in that row. I will construct a concise expression that explicitly states the target grid has an "X" at row 1, column 3, and notes the surrounding "X"s at columns 2 and 4, emphasizing that this specific top{-}row configuration is unique to the target and absent in the distractors.} \\
\\ 
\\ 
\texttt{You are given three grids, where each of them is 5 by 5 in size.} \\
\texttt{Grids have empty cells marked with "▢" and filled cells marked with "X".} \\
\texttt{Your task is to generate a referring expression that best describes the target grid while distinguishing it from the two other distractor grids.} \\
\texttt{The first grid is the target grid, and the following two grids are the distractors.} \\
\\ 
\texttt{Target grid:} \\
\\ 
\texttt{□ □ X X □} \\
\texttt{X □ □ X X} \\
\texttt{□ X X □ X} \\
\texttt{X X □ □ X} \\
\texttt{□ X X □ □} \\
\\ 
\texttt{Distractor grid 1:} \\
\\ 
\texttt{X □ X X □} \\
\texttt{□ X □ X X} \\
\texttt{X X □ □ X} \\
\texttt{□ X X □ X} \\
\texttt{X □ X □ □} \\
\\ 
\texttt{Distractor grid 2:} \\
\\ 
\texttt{□ X X □ X} \\
\texttt{X □ X X □} \\
\texttt{□ X □ X X} \\
\texttt{X □ X □ X} \\
\texttt{□ X X □ □} \\
\\ 
\texttt{Instruction: Describe the target grid.} \\
\texttt{Generate the referring expression starting with the tag "Expression: " for the given target grid. Omit any other text.} \\
            }
        }
    }
    & \\ \\

    \theutterance \stepcounter{utterance}  
    & \multicolumn{4}{p{0.6\linewidth}}{
        \cellcolor[rgb]{0.9,0.9,0.9}{
            \makecell[{{p{\linewidth}}}]{
                \texttt{\tiny{[P1$\rangle$GM]}}
                \texttt{Expression: The grid with an "X" at the top row, third column, and "X"s at the second and fourth columns of the top row, forming a horizontal triplet, while the other two grids have "X"s at the first or second column in the top row.} \\
            }
        }
    }
    & & \\ \\

    \theutterance \stepcounter{utterance}  
    & & & \multicolumn{2}{p{0.3\linewidth}}{
        \cellcolor[rgb]{0.9,0.9,0.9}{
            \makecell[{{p{\linewidth}}}]{
                \texttt{\tiny{[GM$|$GM]}}
                \texttt{Expression: The grid with an "X" at the top row, third column, and "X"s at the second and fourth columns of the top row, forming a horizontal triplet, while the other two grids have "X"s at the first or second column in the top row.} \\
            }
        }
    }
    & & \\ \\

    \theutterance \stepcounter{utterance}  
    & & \multicolumn{4}{p{0.6\linewidth}}{
        \cellcolor[rgb]{0.9,0.9,0.9}{
            \makecell[{{p{\linewidth}}}]{
                \texttt{\tiny{[GM$\rangle$P2]}}
                \texttt{You have attempted this same task before and did not succeed. These are the plans you wrote after those attempts.} \\
\texttt{Trial \#1: New plan: I incorrectly identified the target grid by relying on memory or a hasty glance rather than systematically verifying each condition against the provided grids. Specifically, I failed to check the exact positions of the 'X's in the fifth row (second and third positions) for all options. Next time, I will explicitly list the coordinates of the 'X's for each grid (First, Second, Third) and compare them step{-}by{-}step with the expression's requirements before making a final selection.} \\
\\ 
\texttt{Trial \#2: New plan: I incorrectly identified the target grid by failing to systematically verify the specific spatial constraints of the expression against each grid's layout. Specifically, I missed that the expression requires a vertical column of "X"s in the second column starting from the second row, which is only present in the Third grid (First and Second have "X"s in the first or third rows of that column). Next time, I will explicitly check the top row pattern first to filter options, then verify the vertical column condition in the second column for the remaining candidates before selecting the final answer.} \\
\\ 
\texttt{Trial \#3: New plan: I will systematically verify the expression's two conditions against each grid in order. First, I will check if the top row contains "X"s at columns 3, 4, and 5 (forming a horizontal triplet) and if the second row has "X"s at columns 1 and 5. I will explicitly map the coordinates of every "X" in the top two rows for each grid to ensure no misinterpretation of the pattern before selecting the final answer.} \\
\\ 
\\ 
\texttt{You are given three grids, where each of them is 5 by 5 in size.} \\
\texttt{Grids have empty cells marked with "▢" and filled cells marked with "X".} \\
\texttt{You are also given a referring expression that describes one of the given grids.} \\
\texttt{Your task is to select the grid that matches the given referring expression.} \\
\texttt{Generate only the number (in text) of the grid that matches the given expression by selecting first, second, or third.} \\
\\ 
\texttt{First:} \\
\\ 
\texttt{X □ X X □} \\
\texttt{□ X □ X X} \\
\texttt{X X □ □ X} \\
\texttt{□ X X □ X} \\
\texttt{X □ X □ □} \\
\\ 
\texttt{Second:} \\
\\ 
\texttt{□ □ X X □} \\
\texttt{X □ □ X X} \\
\texttt{□ X X □ X} \\
\texttt{X X □ □ X} \\
\texttt{□ X X □ □} \\
\\ 
\texttt{Third:} \\
\\ 
\texttt{□ X X □ X} \\
\texttt{X □ X X □} \\
\texttt{□ X □ X X} \\
\texttt{X □ X □ X} \\
\texttt{□ X X □ □} \\
\\ 
\texttt{Expression: The grid with an "X" at the top row, third column, and "X"s at the second and fourth columns of the top row, forming a horizontal triplet, while the other two grids have "X"s at the first or second column in the top row.} \\
\texttt{Question: Which grid does the expression refer to?} \\
\texttt{Start with the tag "Answer: ", followed by your selection. Omit any other text.} \\
            }
        }
    }
    & \\ \\

    \theutterance \stepcounter{utterance}  
    & & & \multicolumn{4}{p{0.6\linewidth}}{
        \cellcolor[rgb]{0.9,0.9,0.9}{
            \makecell[{{p{\linewidth}}}]{
                \texttt{\tiny{[GM$\langle$P2]}}
                \texttt{Answer: Third} \\
            }
        }
    }
     \\ \\

    \theutterance \stepcounter{utterance}  
    & & & \multicolumn{2}{p{0.3\linewidth}}{
        \cellcolor[rgb]{0.9,0.9,0.9}{
            \makecell[{{p{\linewidth}}}]{
                \texttt{\tiny{[GM$|$GM]}}
                \texttt{third} \\
            }
        }
    }
    & & \\ \\

    \theutterance \stepcounter{utterance}  
    & \multicolumn{4}{p{0.6\linewidth}}{
        \cellcolor[rgb]{0.95,0.95,0.95}{
            \makecell[{{p{\linewidth}}}]{
                \texttt{\tiny{[P2$\rightarrow$P1]}}
                \texttt{Post{-}episode Reflexion (not sent to GM):} \\
\texttt{New plan: I failed because my previous descriptions were too verbose and failed to explicitly state the single most distinctive feature that separates the target grid from the distractors. The target grid is the only one with an "X" in the top row at the third column (row 1, col 3), whereas the distractors have "X"s at columns 1 or 2 in that row. I will construct a concise expression that explicitly states the target grid has an "X" at row 1, column 3, and notes the surrounding "X"s at columns 2 and 4, emphasizing that this specific top{-}row configuration is unique to the target and absent in the distractors.} \\
            }
        }
    }
    & & \\ \\

    \theutterance \stepcounter{utterance}  
    & & & \multicolumn{4}{p{0.6\linewidth}}{
        \cellcolor[rgb]{0.95,0.95,0.95}{
            \makecell[{{p{\linewidth}}}]{
                \texttt{\tiny{[P1$\rightarrow$P2]}}
                \texttt{Post{-}episode Reflexion (not sent to GM):} \\
\texttt{New plan: I will systematically verify the expression's two conditions against each grid in order. First, I will check if the top row contains "X"s at columns 3, 4, and 5 (forming a horizontal triplet) and if the second row has "X"s at columns 1 and 5. I will explicitly map the coordinates of every "X" in the top two rows for each grid to ensure no misinterpretation of the pattern before selecting the final answer.} \\
            }
        }
    }
     \\ \\

\end{supertabular}
}

\end{document}
